% Options for packages loaded elsewhere
\PassOptionsToPackage{unicode}{hyperref}
\PassOptionsToPackage{hyphens}{url}
%
\documentclass[
]{article}
\usepackage{amsmath,amssymb}
\usepackage{iftex}
\ifPDFTeX
  \usepackage[T1]{fontenc}
  \usepackage[utf8]{inputenc}
  \usepackage{textcomp} % provide euro and other symbols
\else % if luatex or xetex
  \usepackage{unicode-math} % this also loads fontspec
  \defaultfontfeatures{Scale=MatchLowercase}
  \defaultfontfeatures[\rmfamily]{Ligatures=TeX,Scale=1}
\fi
\usepackage{lmodern}
\ifPDFTeX\else
  % xetex/luatex font selection
\fi
% Use upquote if available, for straight quotes in verbatim environments
\IfFileExists{upquote.sty}{\usepackage{upquote}}{}
\IfFileExists{microtype.sty}{% use microtype if available
  \usepackage[]{microtype}
  \UseMicrotypeSet[protrusion]{basicmath} % disable protrusion for tt fonts
}{}
\makeatletter
\@ifundefined{KOMAClassName}{% if non-KOMA class
  \IfFileExists{parskip.sty}{%
    \usepackage{parskip}
  }{% else
    \setlength{\parindent}{0pt}
    \setlength{\parskip}{6pt plus 2pt minus 1pt}}
}{% if KOMA class
  \KOMAoptions{parskip=half}}
\makeatother
\usepackage{xcolor}
\setlength{\emergencystretch}{3em} % prevent overfull lines
\providecommand{\tightlist}{%
  \setlength{\itemsep}{0pt}\setlength{\parskip}{0pt}}
\setcounter{secnumdepth}{-\maxdimen} % remove section numbering
\ifLuaTeX
  \usepackage{selnolig}  % disable illegal ligatures
\fi
\IfFileExists{bookmark.sty}{\usepackage{bookmark}}{\usepackage{hyperref}}
\IfFileExists{xurl.sty}{\usepackage{xurl}}{} % add URL line breaks if available
\urlstyle{same}
\hypersetup{
  hidelinks,
  pdfcreator={LaTeX via pandoc}}

\author{}
\date{}

\begin{document}

\newcommand{\ba}{\begin{aligned}}
\newcommand{\ea}{\end{aligned}}
\newcommand{\bm}{\begin{bmatrix}}
\newcommand{\em}{\end{bmatrix}}
\newcommand{\bc}{\begin{cases}}
\newcommand{\ec}{\end{cases}}
\newcommand{\empty}{\varnothing}
\newcommand{\se}{\subseteq}
\newcommand{\ve}{\varepsilon}
\newcommand{\s}{\sigma}
\newcommand{\0}{\textbf 0}
\newcommand{\1}{\textbf 1}
\newcommand{\I}{\textbf I}
\newcommand{\X}{\textbf X}
\newcommand{\x}{\textbf x}
\newcommand{\y}{\textbf y}
\newcommand{\b}{\textbf b}
\newcommand{\w}{\textbf w}
\newcommand{\P}{\mathcal P}
\newcommand{\S}{\mathcal S}
\newcommand{\T}{\mathcal T}
\newcommand{\A}{\mathcal A}
\newcommand{\B}{\mathcal B}
\newcommand{\C}{\mathcal C}
\newcommand{\L}{\mathcal L}
\newcommand{\c}{\backslash}
\newcommand{\if}{\text{ if }}
\newcommand{\uinf}[1][k]{\bigcup_{{#1} = 1}^{\infty}}
\newcommand{\sinf}[1][k]{\sum_{{#1} = 1}^{\infty}}
\newcommand{\iinf}[1][k]{\bigcap_{{#1} = 1}^{\infty}}
\newcommand{\plim}{\text{plim }}
\newcommand{\var}[2][]{\text{Var}_{#1}\br{#2}}
\newcommand{\e}[2][]{\text{E}_{#1}\br{#2}}
\newcommand{\cov}[3][]{\text{Cov}_{#1}\br{#2,#3}}
\newcommand{\estvar}[1]{\text{Est.Var}\br{#1}}
\newcommand{\estcov}[2]{\text{Est.Cov}\br{#1,#2}}
\newcommand{\asyvar}[1]{\text{Asy.Var}\br{#1}}
\newcommand{\asycov}[2]{\text{Asy.Cov}\br{#1,#2}}
\newcommand{\estasyvar}[1]{\text{Est.Asy.Var}\br{#1}}
\newcommand{\too}{\Longrightarrow}
\newcommand{\sq}{\sqrt}
\newcommand{\dev}[2][i]{({#2}_{#1}-\bar{#2})}
\newcommand{\inv}[1]{{#1}^{-1}}
\newcommand{\add}[1][i]{\sum_{#1 = 1}^n}
\newcommand{\ply}[1][i]{\prod_{#1 = 1}^n}
\newcommand{\br}[1]{\left[#1\right]}
\newcommand{\pr}[1]{\left(#1\right)}
\newcommand{\brace}[1]{\left\{#1\right\}}
\newcommand{\seq}[1][x]{{#1}_1,{#1}_2,\cdots,{#1}_{n}}
\newcommand{\seqadd}[1][x]{{#1}_1+{#1}_2+\cdots+{#1}_{n}}
\newcommand{\dto}{\overset{\text d}{\to}}
\newcommand{\asim}{\overset{\text a}{\sim}}
\newcommand{\st}{\text{ s.t. }}
\newcommand{\om}[1]{\mu^*\pr{#1}}
\newcommand{\lm}[1]{\mu\pr{#1}}
\newcommand{\d}{\text{ d}}
\newcommand{\dm}{\text{ d}\mu}
\newcommand{\dl}{\text{ d}\lambda}
\newcommand{\F}{\mathbb F}
\newcommand{\norm}[1]{\|{#1}\|}
\newcommand{\null}{\text{null }}
\newcommand{\span}{\text{span }}

\hypertarget{metric-spaces}{%
\section{Metric Spaces}\label{metric-spaces}}

\textbf{Definition (Metric)} A metric on a nonempty set \(V\) is a
function \(d:V\times V\to [0,\infty)\) such that

\begin{itemize}
\item
  \(d\pr{f,f} = 0\) for all \(f\in V\).
\item
  If \(f,g\in V\) and \(d\pr{f,g} = 0\), then \(f = g\).
\item
  \(d\pr{f,g} = f\pr{g,f}\).
\item
  (Triangle Inequality) \(d\pr{f,h}\le d\pr{f,g}+d\pr{g,h}\), for all
  \(f,g,h\in V\).
\end{itemize}

\(\pr{V,d}\) is called a metric space.

\textbf{Examples}

\begin{itemize}
\item
  \(V\) is any nonempty set,
  \(d\pr{f,g} = \bc 1\if f\ne g\\0\if f=g\ec\).
\item
  \(V = \R\), \(d\pr{f,g} = \abs{f-g}\).
\item
  (\(l^{\infty}\)-metric) \(V = \R^n\),
  \(d\pr{\pr{x_1,x_2,\cdots,x_n},\pr{y_1,y_2,\cdots,y_n}} = \max\brace{\abs{x_1-y_1},\abs{x_2-y_2}, \cdots,\abs{x_n-y_n}}\).
\item
  \(V = C\pr{\br{0,1}}\), which stands for the collection of all
  continuous functions defined on \(\br{0,1}\).

  \[d\pr{f,g} = \max\brace{\abs{f\pr{t}-g\pr{t}}:t\in\br{0,1}}\]
\item
  \(l^1 = \brace{\pr{a_1,a_2,\cdots}:\sinf\abs{a_k}<\infty}\).

  \[d\pr{\pr{a_1,a_2,\cdots},\pr{b_1,b_2,\cdots}} = \sinf\abs{a_k-b_k}<\infty\]
\end{itemize}

\textbf{Definition (Open Ball)} Let \(\pr{V,d}\) be a metric space,
\(f\in V\), \(r>0\).

\begin{itemize}
\item
  The open ball center center \(f\) radius \(r\) is denoted
  \(B\pr{f,r}\) and defined by

  \[B\pr{f,r} = \brace{g\in V:d\pr{f,g}<r}\]
\item
  The closed ball center center \(f\) radius \(r\) is denoted
  \(\bar B\pr{f,r}\) and defined by

  \[\bar B\pr{f,r} = \brace{g\in V:d\pr{f,g}\le r}\]
\end{itemize}

\textbf{Definition (Open Set)} A subset \(G\se V\) is called open if for
every \(f\in G\), there exits \(r>0\) such that \(B\pr{f,r}\se G\).

\textbf{Definition (Closed Set)} A subset \(F\se V\) is closed if and
only if the complement of \(F\) in \(V\) is open.

\textbf{Definition (Closure)} Let \(E\se V\). The closure of \(E\),
denoted as \(\bar E\), is defined by

\[\bar E:=\brace{g\in V:B\pr{g,\ve}\cap E\ne 0,\forall \ve>0}\]

\textbf{Definition ()} \(\brace{f_k}\se V\), \(f\in V\). We write
\(\lim_{k\to \infty}f_k = k\) to mean
\(\lim_{k\to\infty}d\pr{f,f_k} = 0\).

\emph{\textbf{Proposition}} Let \(\bar E\) be the closure of a set
\(E\). Then

\begin{itemize}
\item
  \(\bar E = \brace{g\in V:\exist \brace{f_k}\se E \st \lim_{k\to\infty} = g}\).
\item
  \(\bar E\) is the intersection of all closed sets containing \(E\).
\item
  \(\bar E\) is closed.
\item
  \(E\) is closed if and only if \(E = \bar E\).
\item
  \(E\) is closed if and only if \(E\) contains the limit of every
  convergent sequence of elements of \(E\).
\end{itemize}

\textbf{Definition (Continuity)} Suppose \(\pr{V,d_v}\) and
\(\pr{W,d_w}\) are metric spaces, and \(T:V\to W\) is a function. For
\(f\in V\), \(T\) is continuous at \(f\) if for every \(\ve>0\), there
exists \(\delta>0\) such that

\[d_w\pr{T\pr{f},T\pr{g}}<\ve\]

given \(d_v\pr{f,g}<\delta\) and \(g\in V\).

\emph{\textbf{Proposition}} All of the followings are equivalent:

\begin{itemize}
\item
  \(T\) is continuous on \(V\).
\item
  If \(\lim_{k\to\infty} f_k = f\), then
  \(\lim_{k\to\infty}T\pr{f_k} = T\pr{f}\).
\item
  \(\inv T\pr{G}\) is open in \(V\) if \(G\) is open in \(W\).
\item
  \(\inv T\pr{F}\) is closed in \(V\) if \(F\) is closed in \(W\).
\end{itemize}

\textbf{Definition (Cauchy Sequence)} A sequence \(f_1,f_2,\cdots\) is
called a Cauchy sequence if for every \(\ve>0\), there exists \(N>0\)
such that

\[d\pr{f_j,f_k}\le \ve\]

for all \(k,j\ge N\).

\emph{\textbf{Proposition}} Every convergence sequence in a metric space
is a Cauchy sequence.

\textbf{Definition} A metric space \(V\) is called complete if every
Cauchy sequence in \(V\) converges to some element in \(V\).

\textbf{Examples}

\begin{itemize}
\item
  \(V = \Q\), \(d\pr{x,y} = \abs{x-y}\). Consider

  \[x_k = \frac{1}{10^{1!}}+\cdots +\frac{1}{10^{k!}}\]

  \(\brace{x_k}\) is Cauchy, but does not converge to a rational number.
\end{itemize}

\emph{\textbf{Proposition}}

\begin{itemize}
\item
  A complete subset of a metric space is closed.
\item
  A closed subset of a metric space is complete.
\end{itemize}

\hypertarget{vector-spaces}{%
\section{Vector Spaces}\label{vector-spaces}}

\textbf{Definition} Suppose \(\pr{X,\S}\) is a measurable space. A
function \(f:X\to \C\) is called \(\S\)-measurable if and only if
\(\Re f:X\to \R\) and \(\Im f:X\to \R\) are \(\S\)-measurable functions.
All measurability results are true for complex-valued functions.

\textbf{Definition} \(f:X\to \C\) such that \(\int \abs{f}<\infty\).
Then \(\int f\dm\) is defined by

\[\int f\dm:=\int \Re f\dm+i\int \Im f\dm\]

\textbf{Definition} Let \(\bar z\) be the complex conjugate of a complex
number \(z\). If \(z =a+bi\), then \(\bar z = a-bi\). Note that

\[\bar{\int f\dm} = \int \bar f\dm \if \int \abs{f}\dm <\infty\]

From now on, \(\mathbb F\) stands for either \(\R\) or \(\C\).

\textbf{Definition (Vector Space)} A vector space (over \(\mathbb F\))
is a set \(V\) with an addition (\(+\)) on \(V\) and a scalar
multiplication (\(\alpha f\), \(\alpha \in \mathbb F\)) on \(V\) that
satisfies

\begin{itemize}
\item
  Commutativity
\item
  Associativity
\item
  Additive identity
\item
  Additive inverse
\item
  Multiplicative identity
\item
  Distributive property
\end{itemize}

\textbf{Definition (\(\F^X\))}

\[\mathbb F^X = \brace{\text{All functions from }X \text{ to }\mathbb F}\]

Functional analysis could be thought of as linear algebra on infinite
dimensional spaces.

\hypertarget{normed-vector-spaces}{%
\subsection{Normed Vector Spaces}\label{normed-vector-spaces}}

\textbf{Definition (Norm)} A norm on a vector space \(V\) (over
\(\mathbb F\)) is a function

\[\|\cdot \|:V\to [0,\infty)\]

such that

\begin{itemize}
\item
  (Positive Definiteness) \(\norm{f} = 0\iff f = 0\),
\item
  (Homogeneity) \(\norm{\alpha f} = \abs{a}\norm{f}\) for all
  \(\alpha \in \F\) and \(f\in V\),
\item
  (Triangle Inequality) \(\norm{f+g}\le \norm{f}+\norm{g}\) for all
  \(f,g\in V\).
\end{itemize}

\textbf{Remarks} Most of the time the first two properties are easy to
check but the hard things come with the proof of Triangle Inequality.

\textbf{Examples}

\begin{itemize}
\item
  \(\pr{\R,\abs{\cdot}}\); \(\pr{\mathbb C,\abs{\cdot}}\).
\item
  \(\F^n = \brace{\pr{a_1,a_2,\cdots,a_n}:a_1,a_2,\cdots,a_n\in \F}\),
  with

  \[\ba
  \norm{\pr{a_1,a_2,\cdots,a_n}}_1 & = \sum_{k = 1}^n\abs{a_k}\\
  \norm{\pr{a_1,a_2,\cdots,a_n}}_2 & = \sq{\sum_{k = 1}^n\abs{a_k}^2}\\
  \norm{\pr{a_1,a_2,\cdots,a_n}}_p & = \pr{\sum_{k = 1}^n\abs{a_k}^p}^{1/p}\\
  \norm{\pr{a_1,a_2,\cdots,a_n}}_{\infty} & = \max\brace{\abs{a_1},\abs{a_2},\cdots,\abs{a_n}} \\
  \ea\]
\item
  \(\ell^1 = \brace{\pr{a_1,a_2,\cdots}:\sinf[n]\abs{a_n}<\infty}\),
  with

  \[\norm{\pr{a_1,a_2,\cdots}} = \sinf[n]a_n\]
\item
  \(X\) is any set. Let \(b\pr{X}\) be the set of all bounded functions.
  \(b\pr{X} = \brace{\text{Bounded functions}: X\to \F}\).

  \[\norm{f} = \sup\brace{\abs{f\pr{x}}:x\in X}\]
\item
  \(C\pr{\br{0,1}} = \brace{\text{Continuou functions}:[0,1)\to\F}\),
  with

  \[\norm{f} = \int_0^1\abs{f}\]
\end{itemize}

A norm will induce a metric on the same vector space. One
straightforward example would be

\[d\pr{f,g} = \norm{f-g}\]

Therefore, any normed vector space is a metric space. However the
converse is false. With this in mind we may apply all stuffs in metric
space to any normed vector space.

We have the meaning of limit:

\[\lim_{k\to\infty}f_k = f\iff \lim_{k\to\infty}\norm{f_k-f} = 0\]

We also have the notion of Cauchy sequence.

\textbf{Definition (Banach Space)} A Banach space is a normed complete
vector space.

\textbf{Examples}

\begin{itemize}
\item
  \(\R,\R^2,\R^n\) are Banach spaces.
\item
  \(C\pr{\br{0,1}} = \brace{\text{Continuou functions}:[0,1)\to\F}\)
  with \(\norm{f} = \max_{\br{0,1}}\abs{f}\) is a Banach space.
\item
  \(\pr{\ell^1,\norm{\cdot}_1}\) is a Banach space.
\item
  \(C\pr{\br{0,1}} = \brace{\text{Continuou functions}:[0,1)\to\F}\)
  with \(\norm{f} = \int_0^1\abs{f}\) is \emph{not} a Banach space.
\item
  \(\pr{\ell^1,\norm{\cdot}_{\infty} = \sup_{k\in \N}\abs{a_k}}\) is
  \emph{not} a Banach space.
\end{itemize}

\textbf{Definition (Series)} Suppose \(g_1,g_2,\cdots\) is a sequence in
a normed vector space \(V\). Then

\[\sinf g_k\text{ exists}\iff \lim_{n\to\infty}\sum_{k = 1}^n\text{ exists}\]

When \(\lim_{n\to\infty}\sum_{k = 1}^n\) exists, the series
\(\sinf g_k\) is said to converge.

\emph{\textbf{Proposition}} \(V\) is a Banach space if and only
\(\sinf g_k\) converges to some element in \(V\) for any sequence
satisfying \(\sinf \norm{g_k}<\infty\).

\emph{\textbf{Proof}}

\begin{itemize}
\item
  \(\too\): Assume \(g_1,g_2,\cdots\) satisfies
  \(\sinf \norm{g_k}<\infty\). Let \(f_j = \sum_{n = 1}^jg_n\). Fix
  \(\ve>0\).

  \[\norm{f_j-f_k} = \norm{\sum_{n = k+1}^jg_n}\le \sum_{n = k+1}^j\norm{g_n}< \ve\]

  which is true for \(k,j\) sufficiently large.
\item
  \(\impliedby\): Suppose \(f_1,f_2,\cdots\) is Cauchy in \(V\). We can
  extract a subsequence of \(\brace{f_j}\) (without relabeling) such
  that

  \[\sinf\norm{f_{k+1}-f_k}<\infty\]

  This implies that \(\sinf \pr{f_{k+1}f_k}\) converges to an element in
  \(V\). Moreover,

  \[\lim_{N\to\infty}\sum_{k = 1}^N\pr{f_{k+1}-f_k} = \lim_{N\to\infty}\pr{f_{N+1}-f_1}\]

  which implies that \(\lim_{N\to\infty}f_{N+1}\) exists and
  \(\lim_{N\to\infty}f_{N+1}\) is also an element in \(V\). The last
  step is that, since the original sequence is Cauchy, and we have
  proved that one of its subsequence is convergent, so is the original
  Cauchy sequence.
\end{itemize}

\hypertarget{bounded-linear-maps}{%
\subsection{Bounded Linear Maps}\label{bounded-linear-maps}}

\textbf{Definition (Linear Map)} Suppose \(V\) and \(W\) are vector
spaces. A function \(T:V\to W\) is called a linear map or linear if

\begin{itemize}
\item
  \(T\pr{f+g} = T\pr{f}+T\pr{g}\) for \(f,g\in V\).
\item
  \(T\pr{\alpha f} = \alpha T\pr{f}\) for \(\alpha\in \F\) and
  \(f\in V\).
\end{itemize}

Basically, a linear map is a mapping that preserves the structure of a
linear space.

\textbf{Remarks} For simplicity we may write \(Tf\) and \(T\pr{f}\)
interchangeably with \(Tf = T\pr{f}\). \(\norm{\cdot}\) is used for both
norms of \(V\) and \(W\), whose meaning depends on the context.

\textbf{Definition (Norm of a Linear Map)} Let \(V\) and \(W\) be normed
vector spaces and \(T:V\to W\) be a linear map. The norm of \(T\),
denoted \(\norm{T}\), is defined by

\[\ba
\norm{T} & = \sup\brace{\norm{T\pr{f}}:f\in V \text{ and }\norm{f}\le 1}\\
& = \sup\brace{\norm{T\pr{f}}:f\in V \text{ and }\norm{f}= 1}\\
& = \sup_{f\ne 0}\brace{\frac{\norm{T\pr{f}}}{\norm{f}}}
\ea\]

\(T\) is called bounded if \(\norm{T}<\infty\). The set of bounded
linear maps from \(V\) to \(W\) is defined by \(\B\pr{V,W}\).

\textbf{Examples}

\begin{itemize}
\item
  \(T:\R^n\to\R^m\) is a linear map. There exists a matrix
  \(A\in \R^{m\times n}\) such that

  \[T\pr{v} = Av\]

  and \(\B\pr{\R^n,\R^m}\simeq \R^{m\times n}\). Let \(\norm{A}\) be the
  largest singular value of \(A\), which is equal to the square root of
  the largest eigenvalue of \(A'A\).
\item
  \(C\pr{[0,3)}\) with \(\norm{f} = \max_{\br{0,3}}\abs{f}\). Consider
  the linear map \(T:C\pr{\br{0,3}}\to C\pr{\br{0,3}}\) so that
  \(Tf:\br{0,3}\to \R\) such that

  \[\pr{Tf}\pr{x} = x^2f\pr{x},\forall x\in [0,3)\]

  Prove that \(\norm{T} = 9\).
\item
  \(V = \brace{\pr{a_1,a_2,\cdots}:a_k = 0 \text{ for all but finitely many }k\in \N}\),
  with \(\norm{\pr{a_1,a_2,\cdots}}_{\infty} = \max\abs{a_k}\). Consider
  the linear map \(T:V\to V\) such that

  \[T\pr{a_1,a_2,a_3,\cdots} = \pr{a_1,2a_2,3a_3,\cdots}\]

  Show that \(\norm{T} = \infty\).
\end{itemize}

\emph{\textbf{Proposition}} Assuming \(V\) is a normed vector space and
\(W\) is a Banach space. \(\B\pr{V,W}\) is a Banach space with
\(\norm{T}\).

\emph{\textbf{Proof}} Consider \(T_1,T_2,\cdots\in \B\pr{{V,W}}\) to be
Cauchy sequence. We want to show that, there is \(T\in \B\pr{V,W}\) such
that \(\lim_{k\to\infty}\norm{T_k-T} = 0\).

Pick any \(f\in V\). Consider the sequence \(T_1f,T_2,\cdots\). We have

\[\norm{T_jf-T_kf}\le \norm{T_j-T_k}\norm{f}\]

which implies that \(T_1f,T_2f,\cdots\) is a Cauchy sequence in \(W\).
Then, there exists \(Tf\in W\) such that

\[\lim_{k\to\infty}T_kf = Tf\in W\]

Define \(T:V\to W\) to be

\[T\pr{f} = \lim_{k\to\infty}T_kf\]

\(T\) is first of all a linear map. Moreover, \(T\) is bounded because

\[\ba
\norm{Tf} & \le \sup_k\brace{\norm{T_kf}}\\
& \le \pr{\sup_{k}\norm{T_k}}\norm{f}<\infty
\ea\]

The last thing to show is that \(\lim_{k\to\infty}\norm{T_k-T} = 0\).
Fix \(k>0\). For any \(\ve>0\), the following is true for \(k\) large
enough:

\[\ba
\norm{\pr{T_k-T}f} & = \lim_{j\to\infty}\norm{\pr{T_k-T_j}f}\\
& \le \lim_{j\to\infty}\norm{\pr{T_k-T_j}}\norm{f}\\
& \le \ve\norm{f}
\ea\]

which implies that

\[\norm{T_k-T}\le \ve \too \lim_{k\to\infty}T_k = T\]

\emph{\textbf{Proposition}} Suppose \(T:V\to W\) is a linear map. \(T\)
is bounded if and only if \(T\) is continuous.

\emph{\textbf{Proof}}

\begin{itemize}
\item
  \(\impliedby\): Assume \(T\) is continuous and not bounded. There
  exists a sequence \(f_1,f_2,\cdots\in V\) such that
  \(\norm{f_k}\le 1\) for all \(k\in \N\) but \(\norm{T_kf}\to\infty\)
  as \(k\to\infty\). Therefore,

  \[\lim_{k\to\infty} \frac{f_k}{\norm{Tf_k}}=0\]

  but

  \[T\pr{\frac{f_k}{\norm{Tf_k}}} = \frac{T\pr{f_k}}{\norm{Tf_k}}\]
\item
  \(\too\): Suppose \(\lim_{k\to\infty}f_k = f\). We have that

  \[\ba
  \norm{Tf_k-Tf} & \le \norm{T}\norm{f_k-f}
  \ea\]

  For the right-hand side,
  \(\lim_{k\to\infty}\norm{T}\norm{f_k-f} = 0\).
\end{itemize}

\hypertarget{bounded-linear-functionals}{%
\subsection{Bounded Linear
Functionals}\label{bounded-linear-functionals}}

\textbf{Definition (Linear Functional)} A \emph{linear functional} is a
linear map from a normed vector space \(V\) to a scalar vector space
(\(\R\) or \(\F\)).

\textbf{Examples}

\begin{itemize}
\item
  \(V = \brace{\pr{a_1,a_2,\cdots}: a_k = 0 \text{ but for finitely many }k\ge 0}\).
  Consider

  \[\phi\pr{a_1,a_2,\cdots} = \sinf a_k\]
\item
  If \(V\) has the norm
  \(\norm{\pr{a_1,a_2,\cdots}} = \sinf \abs{a_k}\), \(\phi\) is bounded.
\item
  If \(V\) has the norm
  \(\norm{\pr{a_1,a_2,\cdots}} = \max_{k\ge 0}\abs{a_k}\), \(\phi\) is
  not bounded.
\end{itemize}

\textbf{Definition} Suppose \(T:V\to W\) is a linear map. The null
space, or kernel, of \(T\), denoted \(\null T\), is defined by

\[\null T = \brace{f\in V:Tf = 0}\]

\textbf{Remarks}: \(\null T\) is a subspace of \(V\). If \(T\) is
continuous, then \(\null T\) is a closed subspace.

\[\null T = \inv T\pr{\brace{0}}\]

\emph{\textbf{Proposition}} Suppose \(\phi:V\to \F\) is a linear
functional that is not identically \(0\). Then the followings are
equivalent:

\begin{itemize}
\item
  \(\phi\) is bounded.
\item
  \(\phi\) is continuous.
\item
  \(\null \phi\) is a closed subspace of \(V\).
\item
  \(\overline{\null \phi}\ne V\).
\end{itemize}

\emph{\textbf{Proof}}

\begin{itemize}
\item
  \(3\too 1\): Suppose \(\null \phi\) is a closed and \(\phi\) is
  unbounded. We hope to reach a contradiction.

  There exits a sequence \(f_1,f_2,\cdots \in V\) such that
  \(\norm{f_k} = 1\) and \(\abs{\phi\pr{f_k}}\ge k\) for each
  \(k\ge 1\). We have

  \[\frac{f_1}{\phi\pr{f_1}}-\frac{f_k}{\phi\pr{f_k}}\in \null \phi\]

  which is only possible if \(\phi\pr{\cdot}\) is a number.

  \[\lim_{k\to\infty}\pr{\frac{f_1}{\phi\pr{f_1}}-\frac{f_k}{\phi\pr{f_k}}} = \frac{f_1}{\phi\pr{f_1}}\]

  But, \(\phi\pr{\frac{f_1}{\phi\pr{f_1}}} = 1\), which is obviously a
  contradiction.
\end{itemize}

Infinite dimensional normed vector space always has a discontinuous or
unbounded linear functional.

\textbf{Definition} Suppose \(\brace{e_k}_{k\in \Gamma}\) is a family of
elements in \(V\).

\begin{itemize}
\item
  \(\brace{e_k}\) is called \emph{linear independent} if there does NOT
  exist a finite nonempty subset \(\Omega\) of \(\Gamma\) and a family
  \(\brace{\alpha_j}_{j\in \Omega}\in \F\c \brace{0}\) such that
  \(\sum_{j\in \Omega}\alpha_je_j = 0\).
\item
  \(\span \brace{e_k} = \brace{\sum_{j\in \Omega}\alpha_je_j:\Omega\se \Gamma, \Omega \text{ is finite}, \alpha_j\in \F}\).
\item
  \(V\) is finite dimensional if and only if there exists a finite set
  \(\Gamma\) and a family \(\brace{e_k}_{k\in \Gamma}\) such that
  \(V = \span \brace{e_k}\). \(V\) is \emph{infinite dimensional} if
  \(V\) is NOT finite dimensional.
\item
  A family \(\brace{e_k}\) in \(V\) is called a \emph{basis} if it is
  linearly independent and \(\span \brace{e_k} = V\).
\end{itemize}

\textbf{Remarks}: By Axiom of choice (or Zoin\textquotesingle s Lemma),
there exists a basis for any vector space, including infinite
dimensional spaces.

\emph{\textbf{Proposition}} If \(V\) is infinite dimensional, then there
exits \(\phi:V\to \F\) such that \(\phi\) is unbounded.

\emph{\textbf{Proof}} Suppose \(\brace{e_k}_{k\in \Gamma}\) is a basis
for \(V\) and \(\Gamma\) is infinite. We can assume by relabeling that
\(\N\se \Gamma\), because each infinite set has at least an countably
infinite one.

Define \(\phi:V\to \F\) by

\[\phi\pr{e_j} = \bc
j\norm{e_j}\ & \if j\in \N\\
0 &\if j\in \Gamma\c \N
\ec\]

By definition,

\[\phi\pr{\sum_{j\in \Omega}a_je_j} = \sum_{j\in \Omega} a_j\phi\pr{e_j}\]

It is trivial to show that \(\norm{\phi} = \infty\).
(\(\norm{\phi}\ge \frac{\abs{\phi\pr{e_j}}}{\norm{e_j}} = j\))

\hypertarget{hahn-banach-theorem}{%
\subsection{Hahn-Banach Theorem}\label{hahn-banach-theorem}}

\textbf{Extension Lemma} Suppose \(V\) is a \emph{real} normed vector
space; \(U\) is a subspace of \(V\); and \(\phi:U\to\R\) is a bounded
linear functional. Suppose \(h\in V\c U\). There exists a bounded linear
functional \(\varphi:U+\R h\to \R\) such that

\[\bc
\varphi_U = \phi\\
\norm{\varphi} = \norm{\phi}
\ec\]

\emph{\textbf{Hahn-Banach Theorem}} Suppose \(V\) is a \emph{real}
normed vector space; \(U\) is a subspace of \(V\); and \(\phi:U\to\R\)
is a bounded linear functional. Suppose \(h\in V\c U\). There exists a
bounded linear functional \(\varphi:V\to \R\) such that

\[\bc
\varphi_U = \phi\\
\norm{\varphi} = \norm{\phi}
\ec\]

Baire\textquotesingle s Category Theorem

\textbf{Definition} Suppose \(U\) is a subset of a metric space \(V\).
The interior of \(U\), denoted \(\text{int }U\), is the set of
\(f\in U\) such that some open ball of \(V\) centered at \(f\) with
positive radius is contained in \(U\).

\textbf{Remark} \(\text{int }U\) is the \emph{largest open set}
contained in \(U\). Note that the closure of \(U\) is the smallest
closed set that contains \(U\).

\textbf{Definition} A subset \(U\) is \emph{dense} in \(V\) if and only
if \(\bar U = V\).

\textbf{Examples} \(\Q\) and \(\R\c \Q\) are both dense in \(\R\).

\textbf{Properties}

\begin{itemize}
\item
  \(U\) is dense in \(V\) if and only if every open subset of \(V\)
  contains at least one element of \(U\).
\item
  \(U\) has empty interior (i.e., \(\text{int }U = \empty\)) if and only
  if \(V\c U\) is dense in \(V\).
\end{itemize}

\emph{\textbf{Baire\textquotesingle s Category Theorem}}

\begin{itemize}
\item
  A complete metric space is \emph{NOT} the \emph{countable} union of
  closed subsets with empty interior.
\item
  The countable intersection of dense open subsets of a complete metric
  space is \emph{NON-empty}.
\end{itemize}

\textbf{Remarks} The intuition of the theorem is stating that, the
countable union of small sets remains small, and the countable
intersection of big sets remains big.

\hypertarget{theorems}{%
\subsection{Theorems}\label{theorems}}

\emph{\textbf{Open Mapping Theorem}} Suppose \(V\) and \(W\) are Banach
spaces. Consider \(T:V\to W\), which is a bounded (or equivalently,
continuous) linear map and subjective (or say, onto). Then \(T\pr{G}\)
is open in \(W\) if \(G\) is open in \(V\).

\textbf{Remarks} If any of the two vector spaces are not Banach, the
theorem goes wrong.

\emph{\textbf{Proof}}

Start with the simplest case. Let \(B\) be the unit ball in \(V\). That
is,

\[B = \B\pr{0,1} = \brace{f\in V:\norm{f}< 1}\]

It suffices to show that \(T\pr{B}\) contains an open ball centered at
\(0\), which is equivalent to prove that \(0\in \text{int }T\pr{B}\).

Indeed, by linearity of \(T\),

\[T\pr{\B\pr{f,a}} = Tf+aT\pr{B}\]

for any \(f\in V\) and \(a>0\).

Let \(G\se V\) be open, and \(f\in G\). By definition of a set being
open, there exists an open ball \(\B\pr{f,a}\se G\). If
\(0\in \text{int }T\pr{B}\),

\[Tf = Tf+0 \in \text{int }T\pr{\B\pr{f,a}}\se \text{int }T\pr{G}\]

By property of interior, \(G\) must be open.

\begin{center}\rule{0.5\linewidth}{0.5pt}\end{center}

The surjectivity and linearity of \(T\) implies that

\[\ba
W = \uinf T\pr{kB} = \uinf kT\pr{B}
\ea\]

Because \(W\) is Banach, \(\bar W = W\). Hence,

\[W = \bar W = \uinf \overline{kT\pr{B}}\]

By Baire\textquotesingle s Theorem, there exists \(k\) such that

\[\text{int }\overline{kT\pr{B}}\ne \empty\]

Because of linearity of \(T\), it follows that
\(\text{int }\overline{T\pr{B}}\ne\empty\).

Then there exists \(f\in B\) such that
\(Tf\in \text{int }\overline{T\pr{B}}\). Hence,

\[0_W\in \overline{T\pr{B-f}}\se \text{int }\overline{T\pr{2B}} = \text{int }\overline{2T\pr{B}}\]

It is implied that there exists \(r>0\) such that

\[\bar{\B}\pr{0_W,2r}\se \text{int }\overline{2T\pr{B}}\]

which immediately implies that

\[\bar{\B}\pr{0_W,r}\se \text{int }\overline{T\pr{B}}\]

\begin{center}\rule{0.5\linewidth}{0.5pt}\end{center}

By definition of closure, take \(h\in W\) and \(\norm{h}\le r\). There
exists \(f\in B\) such that

\[\norm{h-Tf}<\ve\]

Equivalently speaking, for any \(h\in W\), there exists
\(f\in \frac{\norm{h}}{r}B\) such that

\[\norm{h-Tf}<\ve\]

Finally, pick any \(g\in W\) and \(\norm{g}<1\).

\begin{itemize}
\item
  Use \((*)\) with \(h = g\) and \(\ve = \frac12\). There exists
  \(f_1\in \frac1rB\) such that \(\norm{g-Tf_1}<\frac12\).
\item
  Use \((*)\) with \(h = g-Tf_1\) and \(\ve = \frac14\). There exists
  \(f_2\in \frac1{2r}B\) such that \(\norm{g-Tf_1-Tf_2}<\frac14\).
\item
  We continue this process inductively and obtain a sequence
  \(f_1,f_2,\cdots\) such that

  \[\sinf \norm{f_k}\le \sinf \frac{1}{2^{k-1}r} = \frac2r<\infty\]
\end{itemize}

Eventually,

\[\norm{g-Tf_1-Tf_2-\cdots-Tf_n}\le \frac1{2^n}\]

Letting \(n\to \infty\) we have

\[g = \lim_{n\to\infty}T\pr{\sum_{k = 1}^b f_k} = Tf\]

\emph{\textbf{Closed Graph Theorem}} \(T\) is a bounded linear map if
and only if the graph of \(T\) is closed.

\emph{\textbf{Bounded Inverse Theorem}} Let \(T:V\to W\) be a bounded
linear map and bijective. Then \(\inv T\) is bounded and linear.

\emph{\textbf{Principle of Uniform Boundedness}} Suppose \(V\) is
Banach, \(W\) is a normed vector space, and \(\A\) is a family of
bounded linear maps from \(V\) to \(W\) such that

\[\sup\brace{\norm{Tf}:T\in \A}<\infty,\forall f\in V\]

(which is somewhat "pointwise boundedness".)

Then

\[\sup\brace{\norm{T}:T\in \A}<\infty\]

\emph{\textbf{Proof}}

\[V = \uinf[n] \brace{f\in V:\norm{Tf}\le n,\forall T\in \A}:= \uinf[n]V_n\]

Since \(T\) is continuous, \(V_n\) is closed in \(V\), By
Baire\textquotesingle s Theorem, there exists \(V_n\) that has nonempty
interior:

\[\B\pr{h,r}\se V_n,\forall h\in V,r>0\]

Now suppose \(g\in V\) and \(\norm{g}<1\). Thus,

\[rg+h\in \B\pr{h,r}\se V_n\]

which implies that

\[\ba
\norm{Tg} & = \norm{\frac{T\pr{rg+h}}{r}-\frac{Th}{r}}\\
& \le \frac nr+\frac{\norm{h}}{r}\\
\too \norm{T} & \le \frac nr+\frac{\norm{Th}}{r}\\
\too \sup_{T\in \A}\brace{\norm{T}} & \le \frac{n+\sup\brace{\norm{Th}:T\in \A}}{r}<\infty
\ea\]

\end{document}
